
En una rueda de madera de 100 cm de diámetro es necesario colocar un neumático de hierro, cuyo diámetro es 5 mm menor que el de la rueda. ¿En cuántos grados es necesario elevar la temperatura del neumático? El coeficiente de dilatación lineal del hierro es $\alpha_1 = 12 \cdot 10^{-6}$ grad$^{-1}$.